\documentclass[../main]{subfiles}
\begin{document}

\chapter{Número de Reynolds}

\section{Introducción al Número de Reynolds}

El número de Reynolds es un parámetro adimensional\footnote{Cantidad sin unidades físicas asociadas} fundamental en mecánica de fluidos\footnote{Rama de la física que estudia el comportamiento de los fluidos en movimiento o en reposo}. Fue propuesto por Osborne Reynolds en 1883 y se utiliza para predecir el régimen de flujo en diferentes situaciones. Este número relaciona las fuerzas inerciales y las fuerzas viscosas presentes en un fluido, lo que permite caracterizar el movimiento del mismo.

\subsection{Importancia del Número de Reynolds}

El número de Reynolds es crucial para:

\begin{itemize}
    \item Determinar si el flujo es laminar o turbulento.
    \item Predecir la transición entre estos regímenes.
    \item Analizar la similitud dinámica en modelos a escala.
    \item Diseñar sistemas de tuberías y canales.
    \item Optimizar el rendimiento de dispositivos como alas de aviones o hélices.
\end{itemize}

\subsection{Fórmula del Número de Reynolds}

La fórmula general para calcular el número de Reynolds es:

\begin{equation}
    Re = \frac{\rho v L}{\mu}
\end{equation}

Donde:
\begin{itemize}
    \item $Re$: Número de Reynolds (adimensional)
    \item $\rho$: Densidad del fluido (kg/m³)
    \item $v$: Velocidad característica del fluido (m/s)
    \item $L$: Longitud característica (m)
    \item $\mu$: Viscosidad dinámica del fluido (kg/(m·s))
\end{itemize}

Alternativamente, podemos expresar el número de Reynolds utilizando la viscosidad cinemática $\nu$:

\begin{equation}
    Re = \frac{v L}{\nu}
\end{equation}

Donde $\nu = \frac{\mu}{\rho}$ es la viscosidad cinemática (m²/s).

\subsection{Interpretación del Número de Reynolds}

El valor del número de Reynolds nos permite clasificar el régimen de flujo:

\begin{itemize}
    \item $Re < 2300$: Flujo laminar
    \item $2300 < Re < 4000$: Flujo de transición
    \item $Re > 4000$: Flujo turbulento
\end{itemize}

Es importante notar que estos valores son aproximados y pueden variar según la geometría y las condiciones específicas del sistema.

\subsection{Aplicaciones del Número de Reynolds}

El número de Reynolds se aplica en diversos campos de la ingeniería y la ciencia:

\begin{itemize}
    \item Ingeniería hidráulica: Diseño de sistemas de tuberías y canales.
    \item Aerodinámica: Estudio del flujo de aire alrededor de vehículos y estructuras.
    \item Oceanografía: Análisis de corrientes marinas y flujos atmosféricos.
    \item Ingeniería química: Diseño de reactores y equipos de proceso.
    \item Ingeniería biomédica: Estudio del flujo sanguíneo en vasos sanguíneos.
\end{itemize}

\ejemplo{Cálculo del Número de Reynolds}{
Consideremos un flujo de agua a través de una tubería circular. Dados los siguientes datos:

\begin{itemize}
    \item Diámetro de la tubería: $D = 0.05$ m
    \item Velocidad del agua: $v = 2$ m/s
    \item Densidad del agua: $\rho = 1000$ kg/m³
    \item Viscosidad dinámica del agua: $\mu = 0.001$ kg/(m·s)
\end{itemize}

Calculemos el número de Reynolds:

\begin{equation}
    Re = \frac{\rho v D}{\mu} = \frac{1000 \text{ kg/m³} \times 2 \text{ m/s} \times 0.05 \text{ m}}{0.001 \text{ kg/(m·s)}} = 100,000
\end{equation}

Como $Re > 4000$, podemos concluir que el flujo es turbulento.
}

\actividad{Responder el siguiente cuestionario}
\begin{enumerate}
    \item ¿Quién propuso el número de Reynolds y en qué año?
    \item ¿Qué fuerzas relaciona el número de Reynolds?
    \item ¿Cuáles son las variables que intervienen en la fórmula del número de Reynolds?
    \item ¿Qué representa la longitud característica en la fórmula?
    \item ¿Cuál es la diferencia entre viscosidad dinámica y cinemática?
    \item ¿Qué rangos de valores del número de Reynolds corresponden a flujo laminar y turbulento?
    \item ¿Por qué es importante el número de Reynolds en la ingeniería hidráulica?
    \item ¿Cómo se aplica el número de Reynolds en aerodinámica?
    \item ¿Qué papel juega el número de Reynolds en el diseño de modelos a escala?
    \item ¿En qué unidades se expresa el número de Reynolds?
\end{enumerate}

\actividad{Resolver los siguientes problemas}
\begin{enumerate}
    \item Calcular el número de Reynolds para el aire que fluye a 15 m/s a través de un tubo de 10 cm de diámetro. La densidad del aire es 1.225 kg/m³ y su viscosidad dinámica es 1.789 × 10⁻⁵ kg/(m·s).
    \item Determinar la velocidad del agua en una tubería de 2 cm de diámetro si el número de Reynolds es 2000. La viscosidad cinemática del agua es 1 × 10⁻⁶ m²/s.
    \item Un fluido con densidad 800 kg/m³ y viscosidad dinámica 0.01 kg/(m·s) fluye por un canal rectangular. Si la velocidad es 0.5 m/s y el número de Reynolds es 1000, calcular la longitud característica del canal.
    \item ¿Cuál es la viscosidad dinámica de un fluido que fluye a 3 m/s en una tubería de 5 cm de diámetro si su número de Reynolds es 15000 y su densidad es 920 kg/m³?
    \item Calcular la densidad de un fluido que fluye en una tubería de 8 cm de diámetro con una velocidad de 2.5 m/s. El número de Reynolds es 20000 y la viscosidad dinámica del fluido es 0.002 kg/(m·s).
    \item Un líquido fluye a través de una tubería circular con un número de Reynolds de 1800. Si la velocidad se duplica manteniendo las demás condiciones constantes, ¿cuál será el nuevo número de Reynolds?
    \item El agua fluye a través de una tubería de 4 cm de diámetro con un número de Reynolds de 3500. Calcular la velocidad del flujo si la viscosidad cinemática del agua es 1 × 10⁻⁶ m²/s.
    \item Un fluido tiene una viscosidad cinemática de 2 × 10⁻⁵ m²/s. Si fluye a 1.5 m/s en un canal con una longitud característica de 0.1 m, ¿cuál es su número de Reynolds?
    \item Determinar el diámetro de una tubería si el número de Reynolds es 5000, la velocidad del fluido es 2 m/s y su viscosidad cinemática es 1.5 × 10⁻⁵ m²/s.
    \item Un aceite con viscosidad dinámica de 0.08 kg/(m·s) y densidad 900 kg/m³ fluye por una tubería. Si el número de Reynolds es 500 y el diámetro de la tubería es 3 cm, calcular la velocidad del flujo.
\end{enumerate}

\end{document}